% Préambule du rapport : compilation depuis main.tex.
% Les textes de couverture et les couleurs peuvent être modifiés ci-dessous.
\usepackage[T1]{fontenc}
\usepackage[utf8]{inputenc}
\IfFileExists{french.ldf}{%
  \usepackage[french]{babel}
}{%
  \usepackage{babel}
  \babelprovide[main,import]{french}
}
\usepackage{lmodern}
\usepackage{microtype}
\usepackage[a4paper,left=23mm,right=23mm,top=25mm,bottom=21mm]{geometry}
\usepackage{graphicx}
\usepackage{xcolor}
\usepackage{amsmath,amssymb}
\usepackage{booktabs,tabularx,array,longtable}
\usepackage{ragged2e}
\usepackage{caption}
\usepackage{placeins}
\usepackage{titlesec}
\usepackage{tikz}
\usetikzlibrary{arrows.meta,positioning}
\usepackage{hyperref}
\setlength{\parindent}{0pt}

\definecolor{anthracite}{HTML}{252728}
\definecolor{terreCuite}{HTML}{A5553E}
\definecolor{grisTexte}{HTML}{62666A}
\definecolor{grisCredit}{HTML}{7B7E81}
\newcolumntype{L}{>{\raggedright\arraybackslash}X}
\newcommand{\nomChapitre}{CHAPITRE}
\newcommand{\code}[1]{{\ttfamily\small\detokenize{#1}}}
\newcommand{\sourcecalcul}[1]{\par\vspace{2pt}{\footnotesize\color{grisTexte}Source : #1.\par}}
\newcommand{\bilan}[1]{\par\medskip\noindent{\sffamily\bfseries À retenir. }#1\par}
\newcommand{\figurerapport}[4][0.98]{%
  \begin{figure}[htbp]\centering
  \includegraphics[width=#1\textwidth]{figures/#2}
  \caption{#3}\label{#4}
  \sourcecalcul{calculs du projet, données BAAC ; notebooks 04 et 05}
  \end{figure}%
}
\captionsetup{font=small,labelfont=bf,skip=7pt,justification=raggedright,singlelinecheck=false}
\renewcommand{\figurename}{Figure}
\renewcommand{\tablename}{Tableau}
\setlength{\textfloatsep}{14pt plus 3pt minus 2pt}
\setlength{\intextsep}{12pt plus 3pt minus 2pt}
\setlength{\floatsep}{12pt plus 3pt minus 2pt}
\renewcommand{\topfraction}{0.90}
\renewcommand{\bottomfraction}{0.80}
\renewcommand{\textfraction}{0.08}
\renewcommand{\floatpagefraction}{0.72}
\setcounter{topnumber}{3}
\setcounter{bottomnumber}{2}
\titleformat{\chapter}[display]
  {\sffamily\bfseries\color{anthracite}}
  {\color{terreCuite}\fontsize{11}{14}\selectfont \nomChapitre\ \thechapter}
  {8pt}{\fontsize{24}{29}\selectfont}
  [\vspace{3mm}{\color{terreCuite}\rule{18mm}{1.2pt}}]
\titlespacing*{\chapter}{0pt}{0pt}{17pt}
\titleformat{\section}{\sffamily\large\bfseries\color{anthracite}}{\thesection}{0.7em}{}
\titlespacing*{\section}{0pt}{15pt}{7pt}
\titleformat{\subsection}{\sffamily\normalsize\bfseries\color{anthracite}}{\thesubsection}{0.7em}{}
\titlespacing*{\subsection}{0pt}{10pt}{5pt}
\titleformat{\part}[display]{\centering\sffamily\bfseries\color{anthracite}}
  {\color{terreCuite}\Large PARTIE \thepart}{18pt}{\fontsize{27}{34}\selectfont}

\newcommand{\titreRapport}{Accidents de la route :\\comprendre leur gravité\\et la prédire}
\newcommand{\sousTitreRapport}{Analyse de données, étude statistique\\et modélisation prédictive}
\newcommand{\perimetreRapport}{France métropolitaine\enspace\(\cdot\)\enspace Données BAAC 2019-2024}
\newcommand{\auteurRapport}{Papa Birane Mbengue}
\newcommand{\formationRapport}{Élève ingénieur en informatique - Sup Galilée}
\newcommand{\dateRapport}{Octobre 2026}
\newcommand{\imageCouverture}{figures/photo_accident_DNA.jpg}
\newcommand{\creditPhoto}{Source : journal DNA}

\hypersetup{
  pdftitle={Accidents de la route : comprendre leur gravité et la prédire},
  pdfauthor={Papa Birane Mbengue},
  pdfsubject={Données BAAC 2019-2024 - France métropolitaine},
  hidelinks
}


% Figures du volet prédictif : source précisée par notebook.
\newcommand{\figurerapportML}[5][0.98]{%
  \begin{figure}[htbp]\centering
  \includegraphics[width=#1\textwidth]{figures/#2}
  \caption{#3}\label{#4}
  \sourcecalcul{calculs du projet, données BAAC ; notebook #5}
  \end{figure}%
}
